\documentclass[11pt]{article}
\usepackage[margin=2.1cm]{geometry}
\usepackage{booktabs}
\usepackage{amsmath}
\usepackage{enumitem}
\usepackage{xcolor}
\usepackage[T1]{fontenc}
\usepackage{lmodern}

\setlength{\parindent}{0pt}
\setlength{\parskip}{3.6pt}
\newcommand{\code}[1]{\texttt{#1}}
\newcommand{\hr}{\vspace{2pt}\hrule\vspace{6pt}}

\begin{document}

\begin{center}
{\Large\bfseries EE\,362 --- Computer Organization Laboratory}\\[2pt]
{\large\bfseries Final Lab --- Demonstration and Extension}\\[4pt]
Fall 2026 \quad$\cdot$\quad Lab TA: Gandhi Torres \quad$\cdot$\quad Mon 11/30 / Tue 12/1, 3:45--5:45~PM, ECIII 237
\end{center}
\hr

\section*{1. What happens today}

\begin{center}
\begin{tabular}{lll}
\toprule
3:45--4:15 & \textbf{Demonstrations} & $\sim$4 minutes each, two questions from a fixed bank \\
4:15--5:25 & \textbf{Extension work} & one track, chosen below \\
5:25--5:45 & \textbf{Submission and close} & \\
\bottomrule
\end{tabular}
\end{center}

\textbf{Full report due end of finals week.} Today is the last session; the report
is finished afterwards.

\section*{2. The demonstration}

Show your CPU running the program. Four minutes. Then two questions, drawn from
this bank --- the same two for everyone in your section, so it is fair:

\begin{enumerate}[leftmargin=2em,itemsep=2pt]
\item Point at any cycle in your waveform and tell me what is executing.
\item Which component gave you the most trouble, and how did you find the bug?
\item If I asked you to add one instruction, what would you change, and in what
order?
\item Where is your critical path, and what would you do to make this run faster?
\end{enumerate}

\textbf{If your CPU is not running}, come to office hours before today, or take
the reference CPU and demonstrate that. Nobody should reach this session without a
working machine to show --- and using the reference costs you nothing on the
demonstration, only on Lab 12, where the grade already stands.

\section*{3. The extension --- pick exactly one}

\begin{center}
\begin{tabular}{llp{6.2cm}}
\toprule
& track & best for \\
\midrule
\textbf{A} & New instruction & a clean, certain finish \\
\textbf{B} & Performance measurement & people who like measuring things \\
\textbf{C} & Pipeline and hazards & strong students; extends the 10/14 lecture \\
\bottomrule
\end{tabular}
\end{center}

\subsection*{Track A --- add an instruction, end to end}

Add \code{ORI} to your CPU: decoder row, testbench vector, and a supplied test
program that actually runs it.

\code{ORI rt, rs, imm} computes $rt \leftarrow rs \mathbin{\mathrm{OR}}
\mathrm{zeroext}(imm)$. Opcode \code{001101}. It mirrors \code{ADDI} exactly
except for the operation and the extension --- \code{ORI}'s immediate is a
\emph{bit pattern}, so it zero-extends, the same argument as \code{LUI} in Lab 4.

The supplied program uses \textbf{two} \code{ORI}s in a row on purpose: one alone
would pass even if you had implemented it as ``copy the immediate'', which is what
happens when you clone the \code{LUI} row without thinking.

\textbf{The harder variant, if you want it:} add \code{SUB} (\code{SPECIAL},
func \code{100010}) instead. You will discover something the eight-instruction set
hid from you --- \code{SUB} is not commutative, and your datapath's operand
routing was fixed by instructions that were. Work out what your CPU's \code{SUB}
actually computes, and document the deviation from real MIPS. \emph{That
discovery is worth more than the instruction.}

\subsection*{Track B --- measure what the processor did}

Instantiate \code{PerfCounters} in your CPU and measure the \textbf{instruction
mix}: ALU operations, loads, stores, control flow.

Your CPU executes one instruction per cycle, so CPI is 1 and measuring it tells
you nothing. \textbf{The mix is the interesting number}, because it is a property
of the \emph{program} and it lets you answer questions about hardware you have not
built:

\begin{itemize}[leftmargin=1.6em,itemsep=1pt]
\item if loads and stores took 3 cycles instead of 1, what would CPI become?
\item if every control-flow instruction cost 2 extra cycles, what would it become?
\end{itemize}

This program is roughly half control flow and almost no memory. Work out which
hypothetical hurts more, and why --- that is Amdahl's law measured rather than
quoted.

The four classes must be \textbf{mutually exclusive and exhaustive}: the testbench
checks that they sum to the cycle count, and a classification with a hole in it
fails.

\subsection*{Track C --- one pipeline stage, and everything it breaks}

\begin{quote}
\textbf{A working pipelined CPU is NOT required. Breaking it correctly and
explaining precisely why is full credit.}
\end{quote}

That is the actual grading rule, not encouragement. Read it again before you
decide.

Insert one pipeline register between fetch and execute. Make \emph{arithmetic}
instructions work on the supplied straight-line program. Then document every way
the pipeline is now wrong:

\begin{itemize}[leftmargin=1.6em,itemsep=1pt]
\item the \textbf{control hazard} --- the PC is updated in stage 2, so a wrong
instruction is already fetched after every branch. This is why real MIPS has a
branch delay slot;
\item the \textbf{data hazard} --- \code{ADDI r1,r1,-1} followed immediately by
\code{BLEZ r1} reads \code{r1} while it is still being written;
\item the \textbf{structural hazard} --- and why this design does \emph{not} have
one, which is the reason Harvard architectures exist;
\item what the pipeline register holds at reset, and why all-zeros is harmless.
\end{itemize}

For each hazard that is present, give a \textbf{concrete instruction pair from the
Lab 12 program} that triggers it, one named real-world fix, and one sentence on
what that fix costs.

\textbf{No credit} for a pipeline that ``works'' because the test program has no
branches and no adjacent dependencies. That is not a working pipeline; it is an
untested one.

\section*{4. Finishing the report}

The final submission adds, on top of your twelve chapters:

\begin{itemize}[leftmargin=1.6em,itemsep=1pt]
\item an \textbf{executive summary} of the CPU --- one page, what it is and what
it does;
\item a \textbf{resource and timing summary} across all the labs: one table, the
ALM and register counts and Fmax of each component;
\item the \textbf{extension chapter} for your chosen track;
\item a \textbf{reflection} chapter --- \emph{what I would design differently};
\item the \textbf{Response to feedback} section for Checkpoint 4;
\item a complete \textbf{code appendix}.
\end{itemize}

\textbf{Keep this project.} It is a genuine portfolio piece: a processor designed
from gates up, verified at every stage, with the evidence to prove it. It is the
most convincing thing you can put in front of someone who asks whether you can
actually design hardware.

\section*{5. Grading (100 points)}

\begin{center}
\begin{tabular}{lr}
\toprule
Demonstration: the CPU runs, and you can explain it & 30 \\
Extension: the work itself & 30 \\
Extension: the write-up and analysis & 20 \\
Final report completeness (summary, tables, reflection, appendix) & 20 \\
\midrule
Total & 100 \\
\bottomrule
\end{tabular}
\end{center}

For Track C, the 30 points of ``the work itself'' are for the \textbf{hazard
analysis}, not for a functioning pipeline.

\vspace{10pt}
\hr
{\small Three months ago you wrote a four-row truth table for a latch. Every week
since, the same loop: specification, truth table, equations, VHDL, simulation,
check. Nothing was ever added to the method --- you did the same thing thirteen
times at increasing scale, and at the end of it there is a processor. That is
engineering, and it is the actual thing you are taking away from this course.}

\end{document}
